% ===========================================================================
% 1 Introduction (about 4 pages). The paper improves ONE part of the problem: evidence
% retrieval for scientific claims. Everything else supports that improvement.
% Drafted with AI assistance (Claude). Rewrite in your own words before submitting.
% ===========================================================================
\section{Introduction}
\label{sec:introduction}

\subsection{Background and motivation}
\label{sec:motivation}
% TODO: scientific question answering and literature assistants; the pipeline
% question -> search -> passages -> generated answer with citations; a reader can only
% cite what retrieval finds, so retrieval bounds the quality of every answer.

\subsection{The part of the problem this paper improves}
\label{sec:problem}
% TODO: evidence retrieval for a scientific claim: given a claim, rank the abstracts that
% support or refute it. Why this step: it is shared by every assistant, it is measurable
% against known evidence (SciFact), and its cost (reranking) dominates the pipeline.

\subsection{Research questions}
\label{sec:rq}
% TODO:
% RQ1 Can in-domain fine-tuning with weighted fusion improve evidence retrieval over a
%     standard hybrid pipeline with a reranker, at lower cost?
% RQ2 Why does the chosen fusion weight work? (Proposition 1)
% RQ3 Which kind of claim does each retriever help? (word overlap)
% RQ4 Does the improvement hold on a second biomedical dataset?
% RQ5 Where, and why, does the improved retrieval still fail?

\subsection{Contributions: what is new}
\label{sec:novelty}
% TODO, each with its number:
%  1. the improvement: fine-tuned 22M-parameter encoder + weighted RRF, nDCG@10 0.759 vs
%     0.715 for hybrid + reranker (p = 0.003), Recall@10 0.905, 114 vs 329 ms per claim;
%  2. Proposition 1 (Section 3.6): below a BM25 weight of 61/160 ~ 0.38, weighted RRF can
%     only re-order the dense candidates; this turns choosing the weight into a rule;
%  3. word overlap between a claim and its evidence decides which retriever wins;
%  4. the improvement on a second dataset (RQ4);
%  5. a failure analysis (42% of the hardest misses are evidence problems), answer quality
%     and a measured abstention baseline.

\subsection{How this work differs from SciSpace, ResearchRabbit and Consensus}
\label{sec:tools}
% TODO: ResearchRabbit maps citation networks and does not answer questions; Consensus and
% SciSpace answer with citations but are closed and publish no evaluation of their
% retrieval. This paper isolates the retrieval step those assistants depend on, improves
% it, proves why the improvement works, and measures it openly on known evidence; it runs
% offline on a laptop. Comparison table (Table tab:tools, in the chapter PR).

\subsection{Organisation of the paper}
% TODO: one sentence per chapter.
